% =============================================================================
% WEIMING BEAMER THEME - COMPREHENSIVE EXAMPLE
% =============================================================================
% 
% This file serves as the definitive reference guide for the Weiming Beamer 
% theme. It demonstrates every custom environment, semantic markup command, 
% and structural tool provided by the theme.
%
% Compile this file using pdfLaTeX (via the provided Makefile) to ensure 
% optimal font rendering and progress bar calculation.
%
% Note: All Ancient Greek words are transliterated into Latin characters 
% to ensure safe pdfLaTeX compilation without requiring extra font packages.
% =============================================================================

% We use the 'compress' option to keep the Beamer header routing compact,
% which aligns perfectly with Weiming's minimalist philosophy.
\documentclass[10pt, compress]{beamer}

% Load the custom Weiming theme. Make sure beamerthemeweiming.sty is present.
% Choose between: 
%               - light
%		- dark
%		- teal
\usepackage[light]{beamerthemeweiming}

% =============================================================================
% METADATA SETUP
% =============================================================================
% Define the standard presentation metadata. The Weiming \maketitle macro 
% will consume these elements and output a perfectly centered, clean title page.

\title{Ancient Greek and Absurdity}
\subtitle{Aorist Cats and Dual Parrots}
\author{Pau Amaro Seoane}
\date{\today}

\begin{document}

% -----------------------------------------------------------------------------
% 1. TITLE SLIDE
% -----------------------------------------------------------------------------
% The \maketitle command in Weiming is overridden to force a [plain] slide 
% layout, removing all navigation bars, footers, and progress tracks.
\maketitle

% -----------------------------------------------------------------------------
% 2. SECTION DECLARATION
% -----------------------------------------------------------------------------
% In this theme, \section{} silently updates the progress bar mathematics
% and the Table of Contents, but does NOT insert an automatic blank slide.
\section{The Aorist Tense and Confused Felines}

% -----------------------------------------------------------------------------
% 3. STANDARD CONTENT FRAME & SEMANTIC MARKUP
% -----------------------------------------------------------------------------
\begin{frame}{The Punctiliar Cat}
  
  % \idea{} is a semantic wrapper. It applies no visual formatting, but it 
  % clearly distinguishes the setup of your slide in the source code.
  \idea{The Aorist tense (from \textit{a-oristos}, "without horizon") expresses a simple, undefined action in the past.}
  
  \vspace{0.3cm}
  
  % We use the custom \narrow_itemize environment. Standard LaTeX lists waste
  % enormous amounts of vertical space. This environment packs items tightly.
  \begin{narrow_itemize}
    \item \grey{The Imperfect:} The cat \textit{was knocking} the vase over (ongoing).
    \item \grey{The Perfect:} The cat \textit{has knocked} the vase over (and the owner is currently angry).
    \item \grey{The Aorist:} The cat \textit{knocked} the vase over. It happened. Boom.
  \end{narrow_itemize}
  
  \vspace{0.5cm}
  
  \idea{The Sketch Element:}
  
  \begin{narrow_itemize}
    \item A violently confused cat sits next to shattered porcelain.
    \item The cat did not "used to knock it over". It was a singular, absurd event.
  \end{narrow_itemize}

  \vspace{0.5cm}
  
  % Pair \idea{} with \point{} to drive home the slide's takeaway.
  % \point{} automatically applies the custom 'dimred' color and italics.
  \idea{The Grammatical Reality:} \point{The Aorist cat takes no responsibility for the consequences.}
  
  \vspace{1cm}
  
  % Example of using custom gray citation macros at the bottom of a slide.
  % \myrefs creates a dotted line, followed by your tightly spaced citations.
  % \myREF highlights a specific reference in orange (great for self-citations).
  \myrefs{Reference: \myREF{Smyth 1920, \S 1923} | \Citep{Ministry of Silly Felines, 1971}}
\end{frame}

% -----------------------------------------------------------------------------
% 4. CUSTOM TIKZ BOXES
% -----------------------------------------------------------------------------
\begin{frame}{The Aorist Passive Participle}
  
  % Weiming provides pre-configured TikZ boxes for emphasizing warnings, 
  % rules, or theorems without needing to write raw TikZ code every time.
  \begin{center}
    \begin{tikzpicture}
      % TitleBox: A bold, orange, rounded rectangle with white text.
      \node[TitleBox] (mainbox) {The \textit{lytheis} (having-been-loosed) cat.};
      
      % note: An inverted version (white fill, orange border, black text).
      % We use the 'positioning' library (loaded in the sty) to place it.
      \node[note, below=0.5cm of mainbox] {A passive event completed in the past.};
    \end{tikzpicture}
  \end{center}
  
  \vspace{0.5cm}
  
  \begin{center}
    \begin{tikzpicture}
      % box: A light blue container with a red border.
      \node[box] (infobox) {
        \begin{minipage}{0.8\textwidth}
          \centering
          The Confused Cat Sketch:\\
          \vspace{0.2cm}
          "I wish to register a complaint about this feline!\\
          It is a \textit{having-been-confused} (Aorist Passive) cat!"\\
          "No sir, it's just in a state of present continuous bewilderment!"
        \end{minipage}
      };
      % fancytitle: A helper style to add a tab to the top of 'box'.
      \node[fancytitle, right=10pt] at (infobox.north west) {Dialogue Note};
    \end{tikzpicture}
  \end{center}

\end{frame}

% -----------------------------------------------------------------------------
% 5. PLAIN TRANSITION FRAME
% -----------------------------------------------------------------------------
% \plain strips all headers, footers, and progress bars. 
% It sets the background to the dark primary color with white text.
% This is perfect for chapter breaks, dramatic pauses, or title cards.
\plain{Part II: The Dual Number and Deceased Avians}

% -----------------------------------------------------------------------------
% 6. SECTION UPDATE
% -----------------------------------------------------------------------------
\section{The Dual Case}

% -----------------------------------------------------------------------------
% 7. PLOT-ONLY FRAME
% -----------------------------------------------------------------------------
% \plotslide consumes the entire frame to display a figure. 
% It automatically applies a 1.5s fade-in transition (if supported by your PDF 
% viewer) to soften the visual impact of switching to a bright graph.
% We use 'example-image.pdf' here as a placeholder for a graph showing 
% "Parrot Vitality over Time".
\plotslide{example-image.pdf}

% -----------------------------------------------------------------------------
% 8. PLOT WITH COMMENTS FRAME
% -----------------------------------------------------------------------------
% \plotslidecomment pushes the image to the top of the slide and provides 
% precise spacing for up to 4 comments at the bottom.
% The comments are formatted using the ultra-compact 'figcommentlist' environment.
% Empty brackets at the end signify unused comment slots.
\plotslidecomment{example-image.pdf}
  {Ancient Greek features Singular, Plural, and the rare \textbf{Dual} number.}
  {The Dual is used strictly for pairs (e.g., two eyes, two hands, two parrots).}
  {Customer: "I wish to return \textit{to psittako} (these two parrots)."}
  {Shopkeeper: "They are resting! Both of them! A dual state of rest!"}

% -----------------------------------------------------------------------------
% 9. ANOTHER CONTENT FRAME (REPEATING COMMANDS FOR EMPHASIS)
% -----------------------------------------------------------------------------
\begin{frame}{Declining the Dead Parrots}
  
  \idea{In the second declension, the dual endings are famously sparse:}
  
  \begin{narrow_itemize}
    \item Nominative/Accusative/Vocative: \textit{-o} (e.g., \textit{to psittako} - the two parrots).
    \item Genitive/Dative: \textit{-oin} (e.g., \textit{toin psittakoin} - of/to the two parrots).
  \end{narrow_itemize}
  
  \vspace{0.5cm}
  
  \idea{The resulting sketch dialogue:}
  
  \begin{narrow_itemize}
    \item \grey{Customer:} "I nailed \textit{to psittako} (the two parrots) to the perch!"
    \item \grey{Shopkeeper:} "If you hadn't nailed them, \textit{to psittako} would have muscled up to the bars and \textit{voom}!"
    \item \grey{Customer:} "Mate, \textit{to psittako} are deceased! They have ceased to be!"
  \end{narrow_itemize}

  \vspace{0.5cm}
  
  \idea{The final verdict:} \point{The dual parrots have joined the choir invisible!}
  
  \vspace{0.5cm}
  \myrefs{Reference: \myREF{Chapman \& Cleese 1969} | \Cite{Goodwin's Greek Grammar}}
\end{frame}

% -----------------------------------------------------------------------------
% 10. PLAIN TRANSITION FRAME (AGAIN)
% -----------------------------------------------------------------------------
\plain{Part III: The Optative of Wish}

% -----------------------------------------------------------------------------
% 11. MULTIMEDIA FRAME
% -----------------------------------------------------------------------------
\section{The Optative Mood}

% \slidemovie embeds local video content directly into the PDF.
% Arg 1: [options] (optional, e.g., width, height, autoplay)
% Arg 2: the actual video file path (e.g., animation.mp4).
% Arg 3: a static placeholder image (poster) displayed before clicking.
% Arg 4: the caption text displayed beneath the video.
%
% Note: We use 'example-image.pdf' for both the movie and the poster here
% just so this file compiles out-of-the-box on your machine without needing 
% an actual media file present. In a real presentation, Arg 2 would be a .mp4.
\slidemovie{example-image.pdf}{example-image.pdf}{A video of a cat wishing for a fish: "Eithe ho ichthys pareie!" (Would that the fish were here!)}

\end{document}
